\documentclass[11pt, a4paper]{article}

%\usepackage[T1]{fontenc}
%\usepackage{fullpage}

\usepackage[utf8]{inputenc} % comment when using lualatex
\usepackage[italian]{babel} % lingua e a-capo-sillabato
\usepackage{graphicx}
\usepackage[hidelinks]{hyperref} % link di pagina
\usepackage[bottom]{footmisc} % note appiccicate al fondo della pagina
\usepackage{float} % per posizionamento immagini
\usepackage{amsthm}
\usepackage{fancyhdr}
\usepackage[font=small,labelfont=bf]{caption} % small font for caption (and bold Figure word)
\usepackage{tabularx} %tabelle
\usepackage[table]{xcolor} %colore caselle

\pagestyle{fancy}
\fancyhf{}% Clear header/footer
\fancyfoot[C]{\thepage} %add page number
\fancyhead[C]{\footnotesize\textit{Documento:} D5 \hfill SleepCode \hfill \textit{Versione:} 1.0}
\renewcommand{\headrule}{{\color{red!70}\rule{\textwidth}{2pt}}}
\setlength{\headheight}{22pt}

%\pagestyle{myheadings}
%\markright{John Smith\hfill On page styles\hfill}

\renewcommand\UrlFont{\color{blue}\rmfamily}

\theoremstyle{definition}

\newtheorem{funcreq}{RF} %% numerazione dei requisiti funzionali
\newtheorem{nonfuncreq}{RNF} %% requisiti non funzionali
\newtheorem{backend}{BE}
\newtheorem{frontend}{FE}

\title{Report Finale}

\author{Raffaele \textsc{Castagna}\\
Alberto \textsc{Rovesti}\\
Zeno \textsc{Saletti}}

\newcommand{\groupNumber}{G17}

% footnote without number
\newcommand\blfootnote[1]{%
  \begingroup
  \renewcommand\thefootnote{}\footnote{#1}%
  \addtocounter{footnote}{-1}%
  \endgroup
}

% —

% Web address for the project (if any)
% \newcommand{\homepage}{\url{https://www.}}

% data
\date{\today}

\makeatletter{}

% IL PREAMBOLO FINISCE QUI %%%%%%%%%%%%%%%%%%%%%%%%%%%%%%%%%%%%%%%%%%%%%%%%%%%%






\begin{document}

% La pagina di copertina si trova in un file .tex a parte
% NON MODIFICARE QUESTO COMANDO!!!
\begin{titlepage}
\newcommand{\HRule}{\rule{\linewidth}{0.3mm}} % Defines a new command for horizontal lines, change thickness here
\center % Centre everything on the page

%------------------------------------------------
%	Logo
%------------------------------------------------
\includegraphics[width=0.3\textwidth]{materiale/UniTrento_logo_ITA_colore.png}\\[0.5cm]
%------------------------------------------------
%	Headings
%------------------------------------------------
\textsc{\Large Dipartimento di Ingegneria\\e Scienza dell'Informazione}\\[1.5cm]

{\Huge\textbf{Sleep Code}}\\[0.5cm]
\textsc{\large Progetto per il Corso di Ingegneria del Software}\\
\textsc{\large Anno Accademico 2023-2024}\\[0.5cm]

%------------------------------------------------
%	Title
%------------------------------------------------

\HRule\\[0.4cm]
{\huge\bfseries \@title}\\[0.1cm]
\HRule\\[1cm]

\begin{minipage}{\textwidth}
\begin{flushleft}
\textit{Descrizione:} documento di analisi dei requisiti funzionali, non funzionali, front-end e back-end.
\end{flushleft}
\end{minipage}\\[1.5cm]


\begin{minipage}{0.4\textwidth}
\begin{flushleft}
\large
\textit{Numero documento:} D1
\end{flushleft}
\end{minipage}
\begin{minipage}{0.4\textwidth}
\begin{flushright}
\large
\textit{Versione documento:} 2.4
\end{flushright}
\end{minipage}\\[1.5cm]

%------------------------------------------------
%	Author(s)
%------------------------------------------------
\begin{minipage}{0.4\textwidth}
\begin{flushleft}
\large
\textit{Membri del gruppo:}\\
\@author % Your name
\end{flushleft}
\end{minipage}
~
\begin{minipage}{0.4\textwidth}
\begin{flushright}
\large
\textit{Numero gruppo: }
\groupNumber
\end{flushright}
\end{minipage}

% 	If you don't want a supervisor, uncomment the two lines below and comment the code above
% 	{\large\textit{Author(s)}}\\
% 	\@author % Your name

%------------------------------------------------
%	Date
%------------------------------------------------

\vfill\vfill
\textit{Ultima revisione:}
{\@date}

\end{titlepage}

\tableofcontents\blfootnote{\textbf{Consigli utili per la consultazione del testo:} Se il lettore per file \texttt{.pdf} attualmente in uso lo consente, è possible navigare con più semplicità e velocità all'interno di questo documento cliccando sugli elementi dell'indice.}


\newpage
\section*{Scopo del documento}
Questo documento rappresenta il report finale del progetto. L'attenzione è
rivolta all'impegno dedicato alla realizzazione di tutti i
deliverables e del prototipo finale, l'organizzazione del lavoro e la sua
distribuzione tra i membri del gruppo (specificando non solo punti di forza
ma anche la consapevolezza delle criticità incontrate nel percorso), nonché
alla rilevanza delle nozioni relative all'ingegneria del software acquisite,
e applicate  questo progetto, durante il corso e i seminari.


\newpage
\section{Approcci all'ingegneria del software}
Nella presente sezione vengono riassunti brevemente i punti più significativi
dei seminari tenutisi durante il corso, ponendo particolare attenzione ai
metodi e ai principi dell'ingegneria del software emersi e descritti dai
relatori.

\subsection{BlueTensor}
\subsection{Il metodo Kanban}
\subsection{IBM}
\subsection{Meta}
\subsection{U-Hopper}
\subsection{Red Hat}
\subsection{Microsoft}
\subsection{Sistemi Legacy}
\subsection{Gestione del ciclo di vita del software e modernizzazione}
\subsection{APSS e tecnologie per servizi pubblici}

\newpage
\section{Organizzazione del lavoro}
\subsection{Ruoli e attività}
\begin{center}
  \footnotesize
  \begin{tabularx}{\textwidth}{|c||X||X|}
      \hline
      \cellcolor{red!70}Componente del team & \cellcolor{red!70}Ruoli & \cellcolor{red!70}Attività principali\\
      \hline
      Raffaele Castagna & Sviluppatore, architetto, \textit{concept designer} & Ha contribuito in parte a tutti i deliverable, con particolare attenzione al D4. Ha concepito l'idea originale del progetto e la conseguente struttura, anche architetturale, del progetto finale.\\
      \hline
      Alberto Rovesti & Progettista, supporto cordinatore e correttore & Ha svolto un lavoro con particolare interesse nei deliverable D1, D2 e D3 e al D5. Ha revisionato le deliverable una volta finite e a dato supporto al Project nella fase di cordinazione\\
      \hline
      Zeno Saletti & Project leader, redattore, supporto architetturale & In quanto team leader ha organizzato, gestito e strutturato i meeting e ha coordinato gli sforzi di tutti i membri. Ha provveduto alla stesura definitiva dei documenti in LaTex, della cura della veste grafica.\\
      \hline
  \end{tabularx}
\end{center}

\subsection{Distribuzione del carico di lavoro}
\begin{center}
  \footnotesize
  \begin{tabularx}{\textwidth}{|X||c||c||c||c||c||c|}
      \hline
      \cellcolor{red!70}Componente del team & \cellcolor{red!70}D1 & \cellcolor{red!70}D2 & \cellcolor{red!70}D3 & \cellcolor{red!70}D4 & \cellcolor{red!70}D5 & \cellcolor{red!70}Totale\\
      \hline
      Raffaele Castagna & &&&&&\\
      \hline
      Alberto Rovesti & &&&&&\\
      \hline
      Zeno Saletti & &&&&&\\
      \hline
      \cellcolor{red!70}Totale & &&&&&\\
      \hline
  \end{tabularx}
\end{center}

\section{Criticità}
\section{Autovalutazione}
\begin{center}
  \footnotesize
  \begin{tabularx}{\columnwidth}{|X||c|}
      \hline
      \cellcolor{red!70}Componente del team & \cellcolor{red!70}Voto\\
      \hline
      Raffaele Castagna & \\
      \hline
      Alberto Rovesti & \\
      \hline
      Zeno Saletti & \\
      \hline
  \end{tabularx}
\end{center}

\end{document}